\section{Degeneration characters and the klt weight bound}
\label{sec:degeneration}

For a klt generic pair, Section~\ref{sec:blocks} supplies a finite
comparison cover and separates the action of a bounded power of inertia
into actions on its factors. We now pass to semistable models of these
factors. The degree of the necessary curve cover is unrestricted.
The point of this section is to bound the characters on normalized
logarithmic forms; those bounds will cancel the unrestricted ramification
in the weight calculation.

\subsection{Equivariant semistable models}

We first record the form of semistable reduction that we use. A marked
horizontal divisor in this statement may include both a boundary and
exceptional divisors of a generic log resolution.

\begin{lemma}[Marked semistable reduction]\label{lem:marked-semistable}
Let $C_1$ be a smooth projective complex curve with a marked point $c_1$,
and let a finite cyclic group act on $C_1$, fixing $c_1$.
Consider finitely many geometrically integral normal projective varieties
over $\C(C_1)$, each with a compatible semilinear action of this group
and an invariant finite set of horizontal divisor markings.
After a finite extension of curve fields and replacement by compatible
finite-order lifts of the actions, there are equivariant smooth
projective models for which the marked fibre is reduced simple normal
crossings and its union with the horizontal markings is simple normal
crossings. The generic fibres may be log resolutions.
If a given generic fibre is already smooth and has no markings requiring
resolution, it can be left unchanged.
\end{lemma}

\begin{proof}
Take equivariant projective models by closing graphs of the finitely
many translates, and resolve equivariantly, including the reduced
marked fibre and the horizontal markings. Near the marked fibre the
map to the curve is toroidal: in local coordinates a uniformizer is a
monomial in the vertical coordinates, up to a unit. The horizontal
markings are among the remaining coordinates. Root extraction on the
base, normalization, and the projective subdivision theorem of
\cite[Chapter~IV, Section~3]{KKMSD} give the reduced SNC fibre.
The horizontal coordinates remain transverse throughout this
construction.

Here equivariance can be retained in the subdivision step. First
barycentrically subdivide the finite vertical cone complex so that
stabilizers fix each ray of any stabilized cone. They then act trivially
on that cone's lattice. After removing face self-identifications by
further such subdivision, choose compatible projective subdivisions on
orbit representatives and pull them back. The lattices used on the
combinatorial quotient are the original cone lattices, not the
invariant lattices of a geometric quotient. The regular height-one
subdivision theorem after sufficiently divisible ramification applies
to this finite complex. It applies with identical subdivisions to the
copies of strata that split after normalization.

This is the equivariant marked construction used in
\cite[Section~4]{U}; the marked resolution and finite-group extension
are also described in \cite[Paragraphs~16 and~21]{KNX}.
Taking a common further extension handles the finite list of models.
One may take a Galois closure and further roots of an original local
parameter, so that the extension remains Galois over any specified
original curve field. Compatible automorphisms extend to the
compositum; they have finite order. Equivariant resolution away from
the marked fibre gives projective models over the complete curve.
\end{proof}

Apply this lemma to the factors from Corollary~\ref{cor:block-forms}.
Let $C'\to C$ be the resulting Galois curve cover, let $c'$ lie over $c$,
and write
\[
 K'=\C(C'),\qquad l=e(c'/c),\qquad b=s!.
\]
The compatible lift constructed after Corollary~\ref{cor:block-forms}
gives a finite-order automorphism $g$ of the comparison cover whose
base action generates inertia. Its power $\sigma=g^b$ acts separately
on the factors.

Choose a rational uniformizer $u$ at $c'$. Its leading transformation
under $g$ is
\begin{equation}\label{eq:inertia-parameter}
                         g^*u\sim\zeta u,
\end{equation}
where $\zeta$ is a primitive $l$th root of unity.
Only the leading coefficient is used; it is unnecessary to require
$g^*u=\zeta u$ as an equality of rational functions.

For a factor $(R_i,H_i)$, let $p_i=m$ on the rationally connected factor
and $p_i=1$ on the other factors. We have a rational $p_i$-canonical form
$\eta_i$ satisfying
\[
                   \Div(\eta_i)+p_iH_i=0
\]
on its generic normal model. On the semistable resolution choose the
horizontal boundary $\widehat H_i$ to equal the strict transform of
$H_i$ plus the positive parts of the horizontal crepant exceptional
coefficients. These coefficients are less than one. On the abelian
and nonabelian Beauville--Bogomolov factors we have $\widehat H_i=0$,
because their singularities are canonical.

\begin{lemma}[Normalization and products]\label{lem:normalized-product}
Multiplying each $\eta_i$ by an integral power of $u$, one can arrange
$w_u(\eta_i)=0$. For these normalized forms, the exterior product in
common degree $m$ has weight zero. On the comparison cover there is a
base function $a\in K'^*$ such that
\begin{equation}\label{eq:product-weight}
 \phi=a\bigwedge_i\eta_i^{\otimes m/p_i},
 \qquad
 l\,w_z(\phi)=\frac{\ord_{c'}a}{m}.
\end{equation}
Here the exterior product denotes the tensor power of the relative
canonical product identification, with a fixed order of factors.
\end{lemma}

\begin{proof}
On a semistable model every component of the marked fibre has
multiplicity one. The computation of weights on a log-smooth model
in Lemma~\ref{lem:snc-weight} shows that $p_iw_u(\eta_i)$ is an
integer. Multiplication by $u^n$ changes the weight by $n/p_i$, so
normalization is possible. The logarithmic divisor
\[
 \Div\bigl(\eta_i\wedge(du/u)^{\otimes p_i}\bigr)
       +p_i(T_i+\widehat H_i)
\]
is then effective near the marked fibre, and its coefficient on at
least one component of $T_i$ is zero.

We check the product assertion before passing to the comparison
cover. On the fibre product of the semistable models, the local
vertical equations have the form
\[
                   \prod_j x_{ij}=u
                    \quad\text{for each factor }i.
\]
Horizontal coordinates are independent. Include their SNC markings
in the toroidal boundary. Relative logarithmic top forms multiply
over the logarithmic curve with exactly one base differential
$du/u$ in the absolute form. The product therefore acquires no
additional power of $u$. In logarithmic coordinates the resulting
form has at most full logarithmic poles. This remains true on
toroidal resolutions, since their logarithmic canonical generators
pull back to logarithmic generators. Thus its weight is nonnegative.

For equality, choose on each factor a fibre component of zero
logarithmic order and a general point away from all other markings.
The morphism to the curve is smooth there. The product of these
open subsets supplies a fibre component of zero logarithmic order,
so the weight is zero. Finite pullback to the comparison cover
preserves this weight by Lemma~\ref{lem:weight-rules}.

On its normal proper generic fibre, the pulled-back product and
$\phi$ have the same divisor: both trivialize the same degree-$m$
log canonical divisor. Their ratio has zero divisor and hence lies
in the constant field $K'$. The scalar rule and the ramification
rule for weights now give \eqref{eq:product-weight}.
\end{proof}

Since $\sigma$ preserves each factor pair, write
\[
 \sigma^*\eta_i=d_i\eta_i,\qquad d_i\in K'^*.
\]
Both sides have weight zero, hence $\ord_{c'}d_i=0$. Put
$\lambda_i=d_i(c')$. Finite order of $\sigma$ and the fact that it
fixes $c'$ imply that each $\lambda_i$ is a root of unity.
The form $\phi$ descends to the original field and is fixed by $g$.
Taking leading coefficients of its transformation in
\eqref{eq:product-weight} yields
\begin{equation}\label{eq:character-cancellation}
              \zeta^{b\ord_{c'}a}
                       \prod_i\lambda_i^{m/p_i}=1.
\end{equation}
There is no permutation sign in this formula: $\sigma$ preserves
each ordered factor.

We have reduced the weight problem to one concrete question:
can the orders of the $\lambda_i$ be bounded using only $s$ and $m$?
The following subsections treat abelian factors and the remaining
factors separately.

\subsection{Abelian factors}

\begin{lemma}[The abelian character]\label{lem:abelian-character}
For an abelian factor of dimension $h$, the order of its normalized
character $\lambda_i$ is bounded in terms of $h$ alone.
\end{lemma}

\begin{proof}
Here $p_i=1$. On the semistable model the normalized form is a regular
relative logarithmic top form. It is not divisible by $u$ as such a
section, since its logarithmic order is zero on some fibre component.
Consequently it frames the extended top Hodge line. This identification
is the semistable comparison between logarithmic de Rham cohomology
and Deligne's canonical extension, with nilpotent residue
\cite{Deligne,Steenbrink}. The action on the central fibre of this line
is $\lambda_i$.

The integral local system in degree $h$, modulo torsion, has rank
$\binom{2h}{h}$. Indeed its smooth fibres are abelian varieties; our
models were unchanged on the smooth generic fibre. In a small disc
choose an analytic coordinate linearizing the finite base action.
On the universal cover of the punctured disc the action, combined
with parallel transport, is an integral matrix $A$ commuting with
the unipotent monodromy $T$. Since the geometric action has finite
order, a power of $A$ is a power of $T$.

Passing from flat frames to canonical-extension frames changes $A$
by a commuting factor of the form $\exp(c\log T)$, which is unipotent.
The eigenvalues are therefore unchanged. Since the Hodge line is
invariant, $\lambda_i$ is an eigenvalue of the integral matrix $A$.
If its order is $n$, its cyclotomic polynomial divides the
characteristic polynomial of $A$, and
\[
                         \varphi(n)\leq\binom{2h}{h}.
\]
Only finitely many positive integers satisfy this inequality.
Their least common multiple bounds and annihilates the character
orders in the stated dimension.
\end{proof}

\subsection{An equivariant model with a log generator}

For the remaining factors the ordinary Betti numbers are not bounded
by the argument. Instead we use the much smaller structure-sheaf
cohomology and a normal component of a special fibre. This requires
a model on which the normalized logarithmic form generates everywhere
near that fibre.
An exact log trivialization in degree $p_i$ then has a residue in the
same degree on a normal fibre component. The coefficients of its
different must consequently lie in $\{0,1/p_i,\ldots,1\}$, allowing
the normal index theorem to bound the residue character.

\begin{lemma}[Equivariant good model]\label{lem:equivariant-good}
Let $(R_i,H_i)$ be a rationally connected or nonabelian
Beauville--Bogomolov factor of dimension $h>0$ above, with
$p_i\in\{1,m\}$ and normalized form $\eta_i$.
There is a projective equivariant model $Y\to C'$ such that, near $c'$,
\begin{enumerate}
\item $Y$ is klt and its fibre $T=Y_{c'}$ is reduced Cartier;
\item with $J$ the transformed horizontal boundary,
      $(Y,T+J)$ is dlt and $K_Y+J$ is semiample over $C'$;
\item for $\Omega_i=\eta_i\wedge(du/u)^{\otimes p_i}$,
\begin{equation}\label{eq:log-generator}
                     \Div_Y(\Omega_i)+p_i(T+J)=0.
\end{equation}
\end{enumerate}
The geometric generic fibre is birational to $R_i$.
\end{lemma}

\begin{proof}
Start with the equivariant smooth semistable model $Q$ from
Lemma~\ref{lem:marked-semistable}. The pair $(Q,\widehat H_i)$ is klt;
adding the reduced fibre $Q_{c'}$ gives a dlt pair. Since this fibre
is a base pullback, a $(K_Q+\widehat H_i)$-negative program over the
curve is also negative for the adjoint with this fibre added.

We construct the program equivariantly. Let $G$ be the finite cyclic
group generated by $\sigma$ and take the quotient
$\pi:Q\to\overline Q=Q/G$. Define the branch-corrected boundary
$\overline H$ by
\[
          K_Q+\widehat H_i
                 =\pi^*(K_{\overline Q}+\overline H).
\]
Its coefficients are $1-(1-\delta)/e$, where $\delta<1$ is an
upstairs coefficient and $e$ a ramification index. They lie in $[0,1)$,
and finite-map discrepancy comparison makes
$(\overline Q,\overline H)$ klt. A finite quotient of a smooth variety
is $\Q$-factorial.

On the generic fibre upstairs the adjoint has a nonzero section in
a sufficiently divisible degree: the log-resolution error is effective
and exceptional over the log Calabi--Yau factor. Multiplying the
finitely many translates of this section gives an invariant section
in a divisible degree downstairs. The horizontal section inequalities
descend by the finite-map formula. Thus Lemma~\ref{lem:curve-mmp}
applies to the quotient pair over $C'/G$ and gives a terminating
program with relatively semiample output.

Lift each step by normalization in the upstairs function field,
using Stein factorization for contractions. These are the usual
finite-group equivariant MMP diagrams; see also
\cite[Complement~3 and Paragraph~21]{KNX}.
They remain over $C'$, since functions integral over the downstairs
curve extend on the normalizations. Finite pullback preserves relative
negativity and relative ampleness. The canonical pullback equality
continues in codimension one because no step extracts divisors.
In particular the lifted pairs remain klt. The invariant divisor $J$
is rational Cartier: it descends rationally to the $\Q$-factorial
quotient, and its finite pullback is $J$. The displayed canonical
pullback equality makes $K_Y+J$ rational Cartier as well, and their
difference makes $K_Y$ rational Cartier.

For completeness, ordinary $\Q$-factoriality upstairs is unnecessary
for dlt preservation here. The lifted negative contraction or flip
has the same discrepancy comparison as an ordinary MMP step.
Adding the marked fibre changes the adjoint by a pullback and leaves
this comparison unchanged. Discrepancies strictly improve for centres
in the affected locus; hence every lc centre on the output meets the
unchanged locus of an input lc centre. It therefore meets its SNC
locus. This is the generic-SNC characterization of dlt, and gives dlt
on each output; compare \cite[Section~3, before Lemma~16]{dFKX}.

Since no divisors are extracted, every surviving marked fibre
component still has multiplicity one. The fibre is Cartier, and the
klt total space is Cohen--Macaulay. It has no embedded fibre components,
so the fibre is reduced. Relative semiampleness pulls back from the
quotient output.

It remains to prove the exact equality \eqref{eq:log-generator}.
The generic adjoint has Kodaira dimension zero, and the steps preserve
its divisible section spaces. The relative semiample contraction
therefore has zero-dimensional generic image. Connected fibres
identify its Stein image with $C'$, so $K_Y+J\sim_{\Q,C'}0$.
The logarithmic divisor on the left of \eqref{eq:log-generator}
is effective near $c'$ by normalization, and stays effective by
pushforward along the program. Its horizontal part is effective
and rationally trivial on the proper generic fibre, hence zero.
Its remaining vertical part is relatively rationally principal.
A function with vertical divisor is constant on the normal proper
generic fibre, and so comes from $K'$. Thus this divisor is a
rational multiple of $T$ near $c'$.

If the multiple were positive, the dlt discrepancy calculation would
make every weight quotient positive, contradicting
$w_u(\eta_i)=0$. It is therefore zero, proving
\eqref{eq:log-generator}.
\end{proof}

\subsection{Fixed points and normal component characters}

\begin{lemma}[A character on a normal log Calabi--Yau pair]
\label{lem:normal-character}
Fix $h\geq0$ and $p\geq1$. Let $(A,\Delta)$ be a normal integral
projective lc pair over $\C$, and let $\omega$ be a nonzero rational
$p$-canonical form with
\[
              \Delta\geq0,\qquad \Div(\omega)+p\Delta=0.
\]
If a finite-order automorphism $\gamma$ preserves the pair and
$\gamma^*\omega=\lambda\omega$, the order of $\lambda$ divides
an integer depending only on $h=\dim A$ and $p$.
\end{lemma}

\begin{proof}
The divisor of $\omega$ is integral, so the coefficients of $\Delta$
belong to $\{0,1/p,\ldots,1\}$. On the normal quotient
$A/\langle\gamma\rangle$ the crepant branch boundary has coefficients in
\[
 \Psi_p=
 \left\{1-\frac{1-\delta}{e}:
       \delta\in\{0,1/p,\ldots,1\},\ e\in\Z_{>0}\right\}.
\]
This is a rational DCC subset of $[0,1]$. An invariant tensor power
of $\omega$ descends and shows that the quotient log canonical divisor
is rational Cartier and rationally linearly trivial. Finite-map
discrepancy comparison makes the quotient pair lc.

Apply Theorem~\ref{thm:normal-index} in dimension $h$ with coefficient
set $\Psi_p$. Choose its principalizing degree $a$ divisible also by
$p$. The pullback of a degree-$a$ trivialization downstairs is invariant
and has the same divisor as $\omega^{\otimes a/p}$. Their ratio is
constant on the normal proper variety $A$. Consequently
$\lambda^{a/p}=1$. This gives the asserted uniform integer.
\end{proof}

\begin{lemma}[Characters of the other factors]\label{lem:other-characters}
For the normalized rationally connected, Calabi--Yau, and symplectic
factors, the orders of the $\lambda_i$ are bounded in terms of $s$ and $m$.
\end{lemma}

\begin{proof}
Use the model of Lemma~\ref{lem:equivariant-good} and write
$h=\dim R_i$. It is flat over the smooth curve, since its integral
total space is torsion-free over the local discrete valuation rings.
The central fibre $T$ is a union of lc centres of the dlt pair
$(Y,T+J)$ and is therefore Du Bois. Cohomology and base change for a
proper flat family with Du Bois special fibre imply local constancy
of $h^j(\OO)$ near that fibre
\cite[Theorem~2.62 and Corollary~2.64]{KollarFamilies}.
After shrinking, the other fibres are klt, as can be checked on a
resolution. Rational singularities and birational invariance identify
their structure-sheaf cohomology with that of a resolution of the
generic factor.

For the rationally connected factor this gives
\[
                   H^j(T,\OO_T)=0\quad(j>0),\qquad h^0=1.
\]
For a Calabi--Yau factor the only nonzero groups are in degrees $0$
and $h$, each of dimension one. For a symplectic factor there is
one dimension in every even degree and none in odd degree.
These identifications follow from extension of reflexive forms,
the defining form algebras of the factors, and Hodge symmetry on
resolutions \cite{GKKP,HP}.

We use the following consequence of coherent Lefschetz:
a finite-order automorphism of a projective scheme with nonzero
alternating trace on $H^\bullet(\OO)$ has a fixed point.
It applies to this possibly singular fibre. One can either use
\cite{BFQ}, or embed the fibre equivariantly in a smooth projective
space and apply \cite[Corollary~5.5]{Donovan} to its pushed-forward
structure sheaf. If the fibre has no fixed point, the sheaf restricts
to zero along each ambient fixed component, and its Lefschetz
contribution vanishes.

In the rationally connected case the alternating trace of $\sigma$
is $1$, so $\sigma$ has a fixed point. In the symplectic case write
$\mu_1,\ldots,\mu_N$ for the eigenvalues on its nonzero cohomology
groups, where $N=h/2+1$. They are nonzero. Their first $N$ power sums
cannot all vanish: Newton's identities would then give
$\prod_j\mu_j=0$. Thus $\sigma^a$ has a fixed point for some
$1\leq a\leq N$.

In the Calabi--Yau case $p_i=1$ and $J=0$.
Equation~\eqref{eq:log-generator} makes $K_Y$ Cartier near $T$.
The klt Cohen--Macaulay total space, and hence its Cartier fibre,
is Gorenstein there. Adjunction trivializes $\omega_T$ by the
residue of $\Omega_i$. The induced scalar on this generator is
$\lambda_i$: the action on $du/u$ has residue one.
Serre duality therefore gives alternating trace
\[
                            1+(-1)^h\lambda_i^{-1}.
\]
If $\lambda_i$ has order at most two it is already bounded.
Otherwise the trace is nonzero, and $\sigma$ has a fixed point.

In every remaining case a power $\sigma^a$, with $a\leq h+1$, fixes
a point of $T$. At most $h+1$ components of the coefficient-one part
of a dlt boundary on an $(h+1)$-fold meet at one point.
Indeed their common intersection
is a union of lc centres, and at the generic point of such a centre
the pair is SNC. A further power of exponent at most $(h+1)!$
therefore preserves a prime component $A$ through that fixed point.

By dlt adjunction $A$ is normal and carries an effective lc different
$\Delta_A$. The rational $p_i$-pluriresidue $\omega_A$ of $\Omega_i$
satisfies
\[
                      \Div(\omega_A)+p_i\Delta_A=0.
\]
This is an equality of divisors: at the generic smooth point it is
the ordinary residue, and at every prime it follows by taking a
Cartier tensor power in adjunction and dividing the resulting
identity. No change of the residue's degree is required.
Naturality of residue shows that the character of the chosen power
on $\omega_A$ is the corresponding power of $\lambda_i$.

Apply Lemma~\ref{lem:normal-character} with $p_i=1$ or $m$ and
$\dim A=h\leq s$. The power used to stabilize $A$ was bounded in
terms of $s$, so the order of $\lambda_i$ is bounded in terms
of $s$ and $m$. Taking least common multiples supplies a single
annihilating integer for all factors.
\end{proof}

\subsection{Completion of the klt case}

\begin{proposition}[The klt weight bound]\label{prop:klt-weight}
For every $s\geq1$ and $m\geq1$ there is a positive integer
$q_{\mathrm{klt}}(s,m)$ with the following property.
Let $K=\C(C)$ for a smooth projective complex curve, let $z$ be a
uniformizer at $c\in C$, and let $(V,B)$ be a geometrically integral
normal projective geometrically klt $s$-fold pair over $K$.
Assume that $K_V$ is rational Cartier, $B\geq0$, and that a rational
$m$-canonical form $\phi$ satisfies $\Div(\phi)+mB=0$.
Then
\[
                         q_{\mathrm{klt}}(s,m)w_z(\phi)\in\Z.
\]
\end{proposition}

\begin{proof}
Use the product and comparison constructions of
Section~\ref{sec:blocks}, and then the normalized models above.
Lemmas~\ref{lem:abelian-character} and~\ref{lem:other-characters}
give a common integer $Q=Q(s,m)>0$ with $\lambda_i^Q=1$ for every $i$.
Equation~\eqref{eq:character-cancellation} implies
\[
                              l\mid bQ\,\ord_{c'}a .
\]
Together with \eqref{eq:product-weight}, this gives
\[
                 mbQ\,w_z(\phi)=\frac{bQ\,\ord_{c'}a}{l}\in\Z.
\]
Thus $q_{\mathrm{klt}}(s,m)=m\,s!\,Q(s,m)$ works.
The reduced generic presentation in Section~\ref{sec:weights}
is rationally Gorenstein, so this is exactly the klt assertion
needed there. The residue induction therefore proves
Theorem~\ref{thm:weight}.
\end{proof}
